%======================================================================================
\section{How $H$ and $H_\mathrm{loc}$ are computed}\label{app:hloc}

$\mathrm{MSD}(\tau)\propto\tau^{2H}$, so $\log_{10}\mathrm{MSD}=2H\log_{10}\tau+c$. \textbf{$H$} is half the slope of the
least-squares line over all 21 lags (equal weight per lag). \textbf{$H_\mathrm{loc}(\tau_c)$} is half the slope of the line fitted on
the lags with $|\log_{10}\tau-\log_{10}\tau_c|\le0.25$ (a factor 1.78 on each side), equal weight; where fewer than
3 lags lie in the window, the 3 nearest are used. It is a windowed regression, not a pointwise derivative.
$\Delta H$ and $\Delta H_\mathrm{loc}$ of two curves are the differences of the fitted values; for a stratified
average they are the same regression on the averaged $\Delta\log\mathrm{MSD}$ (linear: slope of a difference = difference of slopes).
Bands: the same computation on each bootstrap draw, percentiles lag by lag.

\begin{table}[H]
\centering\footnotesize
\begin{tabular}{@{}r l@{\hspace{14pt}}r l@{\hspace{14pt}}r l@{}}
\toprule
$\tau$ (s) & lags used & $\tau$ (s) & lags used & $\tau$ (s) & lags used \\
\midrule
10 & 10--30 (3 nearest) & 90 & 60--140 & 320 & 210--470 \\
20 & 20--40 (3 nearest) & 100 & 60--170 & 390 & 260--580 \\
30 & 20--50 & 110 & 80--170 & 470 & 320--710 \\
40 & 30--60 & 140 & 80--210 & 580 & 390--1000 \\
50 & 30--80 & 170 & 100--260 & 710 & 470--1000 \\
60 & 40--100 & 210 & 140--320 & 870 & 580--1000 \\
80 & 50--140 & 260 & 170--390 & 1000 & 580--1000 \\
\bottomrule
\end{tabular}
\caption{Lags of the window of each $H_\mathrm{loc}$ value. At 10 and 20~s the window is not centred; at 870 and 1000~s
it is cut at 1000~s and the two values use the same lags, so the end of the curve flattens by construction. The window
resolves nothing narrower than about half a decade.}
\end{table}

\textbf{Phase MSD figures (not in this report)} use a different local exponent: the finite-difference derivative of $\log\mathrm{MSD}$
against $\log\Delta$ on 100 lags, without the factor $1/2$ (about 2 for ballistic motion).

\textbf{Ratios and the glide ratio.} Ratios are sums of numerators over sums of denominators over the runs of a population (a run
goes to the stretch that holds its middle). In a small stratum a glide in rising air has almost no height lost, the glide ratio
diverges and its average over strata is unstable; stratified contrasts use the \emph{height lost per km of glide}, the inverse. Pooled tables keep the glide ratio.

%======================================================================================
\section{Checks of the grouping and of the viewer}\label{app:audit}

\begin{itemize}
  \item \textbf{Code against the methodological note.} Every phase was compared with its code (grid, departure, candidate filter, scan of delays,
        contact, gaps, signed delay, null model, configurations, episodes, membership, blocks): all implemented as written, with four differences of text, \textbf{now
        corrected in red in the note}:
        (1) the phases of the run are Vilpellet's, on 1~Hz flights, not the hidden Markov model of xc-lab, so the heading condition acts only between two such flights;
        (2) headings are in the east and north of the first fix of each flight, not of the fix itself (a rotation of about $3^\circ$ at 300~km, against 45$^\circ$);
        (3) a contact is kept only from the later departure to the earlier landing (at most $\delta$ lost at the ends);
        (4) instants in a gap of the track are labelled solo in the tables.
  \item \textbf{Properties of the tables} (chosen run): every pair interval lasts at least $\tau^*=2285$~s and lies where both flights fly after departure; every configuration has at least
        two members; membership covers each flight from takeoff to landing without holes; on 3\,000 flights the group time equals the union of its pair intervals; on 3\,000 configurations the members are
        connected at the midpoint and no active interval leaves; on 2\,000 episodes every configuration but the first starts when another ends and shares a member; every episode lies in one block. The
        counts of the note (154\,598 and 6\,039 flights; 21.7\% and 6.9\% in a group) are reproduced.
  \item \textbf{Viewer.} The \emph{Group flights} tab reads \texttt{groups/} of its data folder: its seven tables are \textbf{byte-identical} to those of the chosen run; the tracks it draws (\texttt{SSD\_DATA}) have the same fixes,
        times, positions and trimmed starts as those the grouping read (\texttt{SSD\_DISANTE}).
  \item \textbf{Analyses.} Tests on synthetic flights for every step, and on every number of the README of the analyses; the decomposition is checked on a synthetic case where the day raises
        a statistic by 10 and the group by 2: pooled 10, stratified 2.
\end{itemize}

%======================================================================================
\section{Which statistics are drawn}\label{app:redundancy}

Unit: a population on a site-day with at least 25 flights (119 units). A statistic is redundant when the others explain most of it across units ($R^2$ of its linear
regression on them); backward elimination drops the most explained until every $R^2\le0.5$. The four shares of time sum to 100\%, so one of them goes first.

\begin{table}[H]
\centering\footnotesize
\begin{tabular}{@{}l l c@{}}
\toprule
statistic & status & R$^2$ on the five drawn \\
\midrule
time in climb (\%) & drawn & 0.23 \\
searches that end in a climb (\%) & drawn & 0.41 \\
climb rate (m/s) & drawn & 0.20 \\
glide speed (km/h) & drawn & 0.26 \\
glide ratio & drawn & 0.38 \\
time in transition (\%) & dropped, number 1, R$^2$ 1.00 when dropped & 0.83 \\
height gained per climb (m) & dropped, number 2, R$^2$ 0.98 when dropped & 0.91 \\
climbs that start after a search (\%) & dropped, number 3, R$^2$ 0.93 when dropped & 0.77 \\
cross-country speed over 30 min (km/h) & dropped, number 4, R$^2$ 0.93 when dropped & 0.90 \\
time in a search in vain (\%) & dropped, number 5, R$^2$ 0.88 when dropped & 0.71 \\
climbs per hour & dropped, number 6, R$^2$ 0.72 when dropped & 0.71 \\
time in a search that ends in a climb (\%) & dropped, number 7, R$^2$ 0.65 when dropped & 0.65 \\
\bottomrule
\end{tabular}

\caption{Five statistics remain; the others are explained by them.}
\end{table}

%======================================================================================
\section{Pooled curves, hang gliders and sizes}\label{app:curves}

\begin{figure}[H]
\centering
\includegraphics[width=0.78\textwidth]{tamsd.pdf}
\caption{TA-MSD and $H_\mathrm{loc}$ of S, G and A, paragliders and hang gliders (hang gliders: 418 group flights, wide bands at long lags).}
\end{figure}

\begin{figure}[H]
\centering
\includegraphics[width=0.78\textwidth]{size_tamsd.pdf}
\caption{By the size of the configuration (orange, darker with more pilots), S (blue) and, dashed, the same flights alone for each size.}
\end{figure}


%======================================================================================
\section{Single site-days}\label{app:sitedays}

Six site-days chosen by a fixed rule (most sizes with 15 flights or more, then most flights, at most two per takeoff site). Bands resample flights (too narrow: flights of one group are not independent).

\begin{figure}[H]
\centering
\includegraphics[width=0.9\textwidth]{siteday_tamsd.pdf}
\caption{TA-MSD and $H_\mathrm{loc}$ on six site-days by the size of the configuration, with the solo flights of the same site-day where they reach 15 flights.}
\end{figure}

\begin{table}[H]
\centering\footnotesize
\begin{tabular}{@{}l l l l l l l@{}}
\toprule
site-day & solo & 2 & 3 & 4 & 5--6 & 7+ \\
\midrule
Saint Hilaire sud, 2022-03-26 & 0.750 & 0.853 & 0.894 & 0.818 &  & 0.994 \\
Saint Hilaire sud, 2017-03-28 & 0.854 & 0.931 & 0.901 & 0.956 &  & 0.873 \\
M\'eruz, 2019-05-01 &  & 0.926 & 0.939 &  & 0.938 &  \\
M\'eruz, 2016-04-28 &  & 0.894 & 0.958 &  & 0.957 &  \\
Chalvet Ouest, 2016-07-19 &  & 0.951 &  &  & 0.892 &  \\
La Sambuy, 2020-07-25 &  & 0.941 & 0.901 &  &  &  \\
\bottomrule
\end{tabular}

\caption{$H_\mathrm{loc}(1000~\mathrm{s})$ by size on each site-day (empty: fewer than 15 flights). Differences between days exceed those between sizes (the climb rate goes from 0.79 to 2.11~m/s across the six days).}
\end{table}

%======================================================================================
\section{Robustness to the threshold}\label{app:robustness}

With $q=0.2$ contacts need to last only 26.4~min: more and shorter groups, up to one contact in five chance.

\begin{table}[H]
\centering\footnotesize
\begin{tabular}{@{}l l l@{}}
\toprule
 & $q=0.1$, $\tau^*=38.1$~min & $q=0.2$, $\tau^*=26.4$~min \\
\midrule
group flights (TA-MSD) & 33\,435 & 48\,777 \\
$H_\mathrm{loc}(1000~\mathrm{s})$, $\Sl$ & 0.837 & 0.836 \\
$H_\mathrm{loc}(1000~\mathrm{s})$, $\A$ & 0.860 & 0.854 \\
$H_\mathrm{loc}(1000~\mathrm{s})$, $\G$ & 0.893 & 0.892 \\
$\Delta H_\mathrm{loc}(1000~\mathrm{s})$, $\G-\A$ same flight & $+$0.037 [$+$0.035, $+$0.039] & $+$0.032 [$+$0.029, $+$0.035] \\
$\Delta H_\mathrm{loc}(1000~\mathrm{s})$, $\G-\A$ same s-d-h & $+$0.043 [$+$0.040, $+$0.046] & $+$0.042 [$+$0.040, $+$0.044] \\
$\Delta H_\mathrm{loc}(1000~\mathrm{s})$, $\A-\Sl$ same s-d-h & $+$0.003 [$+$0.001, $+$0.006] & $+$0.002 [$-$0.001, $+$0.005] \\
$\Delta H_\mathrm{loc}(100~\mathrm{s})$, $\G-\A$ same flight & $-$0.009 [$-$0.009, $-$0.008] & $-$0.016 [$-$0.017, $-$0.015] \\
$\Delta H_\mathrm{loc}(100~\mathrm{s})$, $\G-\A$ same s-d-h & $+$0.002 [$+$0.001, $+$0.002] & $-$0.004 [$-$0.005, $-$0.003] \\
$\Delta$ searches to climb, $\G-\A$ same s-d-h (points) & $+$2.6 [$+$2.3, $+$3.0] & $+$3.4 [$+$3.0, $+$3.7] \\
$\Delta$ climb rate, $\G-\A$ same s-d-h (m/s) & $+$0.029 [$+$0.025, $+$0.034] & $+$0.024 [$+$0.020, $+$0.029] \\
$\Delta$ height lost per km, $\G-\A$ same s-d-h (m/km) & $-$9.5 [$-$10.5, $-$8.6] & $-$11.5 [$-$12.2, $-$10.9] \\
$\Delta$ cross-country speed, $\G-\A$ same s-d-h (km/h) & $+$1.29 [$+$1.16, $+$1.41] & $+$1.14 [$+$1.01, $+$1.30] \\
\bottomrule
\end{tabular}

\caption{The main results with the two thresholds (s-d-h: same site-day and hour). Every conclusion keeps its sign; at 100~s both estimates of the group effect are slightly negative with $q=0.2$: no robust effect within one phase.}
\end{table}

%======================================================================================
\section{Leader minus follower}\label{app:leaders}

\begin{figure}[H]
\centering
\includegraphics[width=0.95\textwidth]{mechanisms_leaders.pdf}
\caption{Leader minus follower, same configuration of two pilots.}
\end{figure}

\begin{table}[H]
\centering\footnotesize
\begin{tabular}{@{}l l@{}}
\toprule
$\Delta$ & leader $-$ follower \\
\midrule
$H_\mathrm{loc}(100~\mathrm{s})$ & $-$0.002 [$-$0.002, $-$0.001] \\
$H_\mathrm{loc}(1000~\mathrm{s})$ & $-$0.003 [$-$0.004, $-$0.001] \\
time in climb (\%) & $+$0.91 [$+$0.57, $+$1.25] \\
searches ending in a climb (\%) & $-$0.01 [$-$1.03, $+$0.83] \\
climb rate (m/s) & $+$0.021 [$+$0.014, $+$0.028] \\
glide speed (km/h) & $+$0.86 [$+$0.74, $+$0.99] \\
height lost per km of glide (m/km) & $+$1.6 [$-$0.9, $+$4.5] \\
cross-country speed (km/h) & $+$0.14 [$+$0.11, $+$0.16] \\
\bottomrule
\end{tabular}

\caption{Leader minus follower, same configuration of two pilots (difference with 5--95\% band).}
\end{table}

\begin{table}[H]
\centering\footnotesize
\begin{tabular}{@{}l l l l@{}}
\toprule
 & climb (\%) & search (\%) & transition (\%) \\
\midrule
entering a group & 28.1 [27.6, 28.5] & 10.3 [10.0, 10.6] & 61.5 [61.0, 61.9] \\
leaving a group & 28.0 [27.6, 28.3] & 11.2 [10.9, 11.5] & 60.6 [60.1, 61.0] \\
share of time, alone & 29.1 [28.9, 29.3] & 10.1 [10.0, 10.2] & 60.8 [60.6, 61.0] \\
share of time, in a group & 33.0 [32.8, 33.2] & 9.4 [9.2, 9.6] & 57.6 [57.3, 57.9] \\
\bottomrule
\end{tabular}

\caption{Phase at the instant of entering or leaving a group, and shares of time in each phase (percent, with 5--95\% bands).}
\end{table}

%======================================================================================
\section{Reproducing the report}\label{app:reproduce}

From \texttt{xc-lab/analyses/groups}, SSD\_DISANTE mounted:
\begin{verbatim}
uv run python tamsd.py; uv run python phase_msd.py; uv run python survival.py
uv run python budget.py; uv run python performance.py; uv run python blocks.py
uv run python redundancy.py; uv run python size.py; uv run python decompose.py
uv run python zones.py; uv run python mechanisms.py; uv run python report/facts.py
uv run python report/tables.py; cd report; latexmk -pdf group_vs_solo.tex
\end{verbatim}
Heavy pieces stay on SSD\_DISANTE, in the \texttt{analyses/} folder of each run; figures and JSON in \texttt{out/<run>/}.
